\answer{ex:05:1}{$\ell\approx17\um$；$\approx100$~天：向全脑广播靠的是导线分支，扩散只是把每个释放位点模糊到$\sim20\um$。}
\answer{ex:05:2}{$r=ab/(1+G)$，$S=1/(1+G)$精确成立；$+1.7\,\%$，而不是$+20\,\%$。}
\answer{ex:05:3}{$G=2$时$T^*=1.21\,\text{s}$，$G=9$时$0.19\,\text{s}$：实际上$G\sim2$--$4$，抑制$3$--$5$倍。}
\answer{ex:05:4}{总频率为$\lambda$时$\mathrm{CV}=1/\sqrt{2\lambda\tau_c}\approx9\cdot10^{-4}$；单个神经元需要$\tau_c=12.5\,\text{s}$。}
\answer{ex:05:5}{两种情况下都是$c\approx0.40$；$\mathrm{CV}=0.050$对$0.158$，相差$\sqrt{10}$倍。各神经元的频率仍与$a_i$成正比：受调节的只是总和。}
\answer{ex:05:6}{$N_1=\overline{A\vee B}$，$N_2=\overline{A\vee N_1}$，$N_3=\overline{B\vee N_1}$，$N_4=\overline{N_2\vee N_3}$，XOR $=N_5=\overline{N_4}$；每次切换用时$\tau_c\ln2$，经过$4$个门的路径使一次XOR需要$2.8\,\text{s}$。}
\answer{ex:05:7}{(a)~$\tau_\gamma=6/\ln3\approx5.5\,\text{s}$。(b)~$\bar D<4\,\text{s}$，而固定延迟为$D<2.2\,\text{s}$：延迟不可预测反而使动物更有耐心。}
\answer{ex:05:8}{$k_2=1/\tau_d=10\,\text{s}^{-1}$，$k_1=1/\tau_d-1/\tau_\gamma=9.5\,\text{s}^{-1}$；$\tau_\gamma=1/(k_2-mk_1)$：$m$变化$+2.6\,\%$时视野加倍（变化$+5.3\,\%$时视野变为无穷大）。}
